\section{Most failures are a false ``done''}
\label{sec:diag}

\paragraph{Taxonomy.} Each episode gets one primary category from fixed rules over its actions, screenshots and final score. \emph{False success}: the agent calls \texttt{answer} on a feasible task scored 0, or claims completion on an infeasible task. \emph{Loop}: at least three consecutive identical actions while less than 1\% of pixels change. \emph{Step limit}: the 50-step budget is exhausted without a loop. At step level, \fsp{} marks the final \texttt{answer} of a false success and \fsm{} that of a true one. A correctly worded refusal of an infeasible task is excluded from both. The rule file was hashed before the official direction analysis (current hash \texttt{\oswRulesSha}); later changes are logged with their hash. One of them, made before any OSWorld AUROC, altered a label: an infeasible task answered with a plainly worded refusal (\texttt{gimp/38f48d40}) moved from \fsp{} to excluded. % source: JOURNAL.md "Labels figés avant toute AUROC OSWorld" and "Règle de label ajoutée avant toute AUROC"
Before the official direction analysis (exploratory directions had been fitted on smaller label sets) we also require that the target category cover at least 25\% of failures and provide at least 40 positive steps (the GO rule).

\paragraph{Result.} On the diagnostic split (\oswNTasks{} tasks, \oswNScored{} scored), Holo-3.1-4B scores \oswMeanScorePct\% (the official OSWorld score, partial credit included, the convention used for every score in this paper) and fully solves \oswNSuccess{} tasks: \oswSuccessPct\%, 95\% Wilson CI \oswSuccessLo--\oswSuccessHi. An independent reproduction with the same public harness reports about 45.5\% at 100 steps on a different task set \citep{realtimegui}; we do not claim the two numbers are comparable. % source: JOURNAL.md "Baseline OSWorld 4B" (RealtimeGUI ~45.5 %)
Of the \oswNFail{} failures (Figure~\ref{fig:osw}), \oswCatFalseSuccess{} (\oswFsPct\%) are false successes, \oswCatStepLimit{} (\oswStepPct\%) have the step limit as primary category (\oswStepAny, or \oswStepAnyPct\%, exhaust the 50 steps at all, loops included), \oswCatLoop{} (\oswLoopPct\%) loop, and \oswCatOther{} is unclassified; \oswNUnscored{} further task is unscored (the ``Other + unsc.'' column of Table~\ref{tab:osw-domains} counts both). \oswFsInfeasible{} false successes are claims of completion on infeasible tasks, and \oswFsFeasible{} are feasible tasks done wrongly or partially. On an infeasible task the correct action is \texttt{answer} with a refusal, so $P(\texttt{answer})$ is the wrong metric there; one such context, Chrome \texttt{480bcfea} at step 4, enters several screens below, and Appendix~\ref{app:infeasible} gives every count with and without it. These are mostly not misworded refusals: among the \infAllowed{} tasks open to construction, \infTasks{} are infeasible (\infEpisodes{} episodes), and of the \infAnswered{} final answers \infClaim{} claim success, while only \infInabilityTasks{} task states its inability, in words the harness rule misses. % source: data/pod_astra/infeas/reviewed_inventory.json, report.md
The category is wider than explicit claims of completion. Read strictly, three of the \oswFsInfeasible{} infeasible false successes of the diagnostic split do not claim completion: two decline the task in words the harness rule misses (\texttt{gimp/38f48d40}, \texttt{vs\_code/7c4cc09e}) and one states that it cannot finish and asks for the user's credentials (\texttt{writer/bb8ccc78}). Of the \oswFsFeasible{} feasible ones, 19 final answers claim completion and one (\texttt{multi\_apps/b337d106}) presents a tutorial instead of applying the setting it describes, on our reading, done by an LLM agent and not validated by humans. The strict count of explicit false claims of completion is therefore 25 of \oswNFail{} failures (49\%); of the other four, two are correct refusals the rule misses and two are answers given instead of finishing, like the episode of Fig.~\ref{fig:teaser}a. % source: final answers in ~/osw/results_diag*/.../traj.jsonl of the 29 false_success tasks listed in data/local/diagnostic_osw_diag.json (readable_decline flags for gimp/38f48d40 and vs_code/7c4cc09e)
They are not early quits: the median false success ends after \oswFsMedianSteps{} steps. Every domain has false successes except Thunderbird, and Chrome, Calc, GIMP and VLC have four or more each (Fig.~\ref{fig:osw}a). % source: generated/tab_osw_domains.tex
False success passes the share criterion of the GO rule, but its \oswLabFsPos{} \fsp{} steps alone do not pass the count criterion. With a second pass over the same tasks, stopped at \rtmPassTwoScored{} scored tasks when the GPU returned to production, \bAdmFsPos{} \fsp{} and \bAdmFsNeg{} \fsm{} steps are admissible, above the required 40 (as of \bStatusUtc{} UTC).

\placefig{fig-osw}

\subsection{The model says ``done'' on most false completions too}
\label{sec:think}
We re-sampled the original model on six screening contexts from five tasks, two per stratum ($k=4$, production sampling, thinking enabled, invalid outputs counted; Fig.~\ref{fig:teaser}c and Table~\ref{tab:causal}, first row). It calls \texttt{answer} on \bPilotFsAnswer/\bPilotFsN{} false-completion samples and \bPilotTrueAnswer/\bPilotTrueN{} true ones, and never on mid-episode contexts (\bPilotHealthyAnswer/\bPilotHealthyN). It reproduces the terminal versus non-terminal distinction exactly, and on these contexts shows no false versus true distinction; on the larger, strictly feasible cohort of \S\ref{sec:levers}, it answers on \pwOrigFspPct\% of false and \pwOrigFsmPct\% of true completions. Any repair must therefore change a decision the model takes in most false and nearly all true completions, and only in the false ones.

The reasoning has the same shape in both cases. On \thN{} re-sampled reasoning traces at completion steps of \thNTasks{} construction tasks, \thFsPosN{} on false and \thFsNegN{} on true completions, every trace opens with a success claim (``The task has been completed successfully'', ``The document has been saved successfully'') at generated-token index \thTokMin--\thTokMax, counted from zero, with no verification word before it (\thNVerif/\thN), then lists the steps taken (median \thMedLen{} tokens). % source: data/pod_astra/think/direction_summary.json (verification_window empty for all events)
Because true completions show it too, this is a property of Holo's output format at completion steps, not a marker of false success. It locates the written conclusion, not the moment the decision forms inside the model. A direction fitted on these tokens separates the two pairs perfectly, but so do context length and a single lexical cue, so we give it no weight; steering along it is negative (Table~\ref{tab:causal}). What the format skips is the check: in the Spotify episode, the reasoning states that the installation cannot finish, and the same step calls \texttt{answer}, reporting the installation as ``initiated successfully'' and handing the password step to the user, which ends the episode unfinished (Fig.~\ref{fig:teaser}a).
